% ---- pg-041 （书页 29） ----

\mvbold{1.5} 利用练习 1.4 中的数据

(a) 画出 $(x_2, x_1)$ 和 $(x_1, x_2)$ 的散布图和边缘点图. 解释这些图形.

(b) 计算 $(x_1, x_2, x_3)$ 的阵列 $\bar{\bm{x}}, \bm{S}$ 和 $\bm{R}$.

\mvbold{1.6} \hyperref[tab:1.5]{表 1.5} 中的数据是对洛杉矶地区空气污染变量的 42 个测量值的记录，这些测量值是在不同日子的中午 12 点采集的（可参见 www. prenhall. com/statistics 的空气污染数据）.

\begin{mvtable}
\phantomsection\label{tab:1.5}\mvtabcap{表 1.5\quad 空气污染数据}
\begin{longtable}{rrrrrrrr}\toprule
风 $(x_1)$ & 日光辐射 $(x_2)$ & $\mathrm{CO}(x_3)$ & $\mathrm{NO}(x_4)$ & $\mathrm{NO}_2(x_5)$ & $\mathrm{O}_3(x_6)$ & $\mathrm{HC}(x_7)$ \\
\midrule
\endfirsthead
\multicolumn{7}{c}{\zihao{6}续表}\\
\toprule
风 $(x_1)$ & 日光辐射 $(x_2)$ & $\mathrm{CO}(x_3)$ & $\mathrm{NO}(x_4)$ & $\mathrm{NO}_2(x_5)$ & $\mathrm{O}_3(x_6)$ & $\mathrm{HC}(x_7)$ \\
\midrule
\endhead
\bottomrule
\endlastfoot

8 & 98 & 7 & 2 & 12 & 8 & 2 \\
7 & 107 & 4 & 3 & 9 & 5 & 3 \\
7 & 103 & 4 & 3 & 5 & 6 & 3 \\
10 & 88 & 5 & 2 & 8 & 15 & 4 \\
6 & 91 & 4 & 2 & 8 & 10 & 3 \\
8 & 90 & 5 & 2 & 12 & 12 & 4 \\
9 & 84 & 7 & 4 & 12 & 15 & 5 \\
5 & 72 & 6 & 4 & 21 & 14 & 4 \\
7 & 82 & 5 & 1 & 11 & 11 & 3 \\
8 & 64 & 5 & 2 & 13 & 9 & 4 \\
6 & 71 & 5 & 4 & 10 & 3 & 3 \\
6 & 91 & 4 & 2 & 12 & 7 & 3 \\
7 & 72 & 7 & 4 & 18 & 10 & 3 \\
10 & 70 & 4 & 2 & 11 & 7 & 3 \\
10 & 72 & 4 & 1 & 8 & 10 & 3 \\
9 & 77 & 4 & 1 & 9 & 10 & 3 \\
8 & 76 & 4 & 1 & 7 & 7 & 3 \\
8 & 71 & 5 & 3 & 16 & 4 & 4 \\
9 & 67 & 4 & 2 & 13 & 2 & 3 \\
9 & 69 & 3 & 3 & 9 & 5 & 3 \\
10 & 62 & 5 & 3 & 14 & 4 & 3 \\
9 & 88 & 4 & 2 & 7 & 6 & 3 \\
8 & 80 & 4 & 2 & 13 & 11 & 4 \\
5 & 30 & 3 & 3 & 5 & 2 & 3 \\
6 & 83 & 5 & 1 & 10 & 23 & 4 \\
8 & 84 & 3 & 2 & 7 & 6 & 3 \\
6 & 78 & 4 & 2 & 11 & 11 & 3 \\
8 & 79 & 2 & 1 & 7 & 10 & 3 \\
6 & 62 & 4 & 3 & 9 & 8 & 3 \\
10 & 37 & 3 & 1 & 7 & 2 & 3 \\
8 & 71 & 4 & 1 & 10 & 7 & 3 \\
7 & 52 & 4 & 1 & 12 & 8 & 4 \\
5 & 48 & 6 & 5 & 8 & 4 & 3 \\
6 & 75 & 4 & 1 & 10 & 24 & 3 \\
10 & 35 & 4 & 1 & 6 & 9 & 2 \\
8 & 85 & 4 & 1 & 9 & 10 & 2 \\
5 & 86 & 3 & 1 & 6 & 12 & 2 \\
5 & 86 & 7 & 2 & 13 & 18 & 3 \\
7 & 79 & 7 & 4 & 9 & 25 & 2 \\
7 & 79 & 5 & 2 & 8 & 6 & 2 \\
6 & 68 & 6 & 2 & 11 & 14 & 3 \\
8 & 40 & 4 & 3 & 6 & 5 & 2 \\
\end{longtable}
\end{mvtable}

{\zihao{6} 资料来源：数据经 G. C. Tiao 教授许可使用.}
